%% =====================================================================
%%  Capítulo 2 -- Sistemas unidimensionales: el flujo sobre la recta
%%  Formato nuevo (noviembre de 2026). Integra: Handout2.tex (2022), el
%%  borrador revisado LibroSisNL2026/capitulos/cap02_unidimensionales.tex
%%  (octubre de 2026) con sus scripts MATLAB/cap02, y las diapositivas
%%  SisUniDim2026.pptx. Agrega la logística con cosecha y la histéresis
%%  del diodo túnel (adelantos de la silla-nodo del capítulo 6).
%%  Correcciones respecto a las versiones anteriores: ver \S\ref{sec:un-cierre}.
%%  Números verificados con herramientas/cap02_numeros.py (42/42) y
%%  codigo/cap02_unidimensionales/verificarCap02.m (37/37 en Octave 8.4).
%% =====================================================================

\hyphenation{sobre unidimensionales}
\capitulo[cap:unidimensionales]{Sistemas unidimensionales: el flujo sobre la recta}
  {La ciencia no es discípula del genio, sino esclava del mundo.}
  {Javier Sampedro}

\begin{objetivos}[partes/02-unidimensionales.html]
  \item Leer el comportamiento de todas las soluciones de $\dot x=f(x)$ en la gráfica de $f$: equilibrios, sentido del flujo y límites, sin resolver la ecuación.
  \item Clasificar equilibrios con la derivada $f'(\eq)$ y reconocer cuándo ese criterio no decide nada.
  \item Explicar por qué un sistema escalar invariante en el tiempo no oscila, y qué hipótesis hay detrás de esa afirmación.
  \item Usar el potencial $V=-\int f$ como un primer ejemplo de función de Lyapunov.
  \item Analizar modelos de varios dominios (población con cosecha, diodo túnel, juego halcón--paloma) y reconocer biestabilidad, regiones de atracción e histéresis.
  \item Reconocer la no unicidad y el escape en tiempo finito, y simularlos sin engañarse.
\end{objetivos}

%% ---------------------------------------------------------------------
\section{¿Por qué empezar por una dimensión?}
\label{sec:un-porque}
%% ---------------------------------------------------------------------

Este capítulo estudia el caso más sencillo posible: un sistema invariante en el tiempo, sin entrada, con una sola variable de estado,
\begin{equation}
  \dot x=f(x),\qquad x(0)=x_0,\qquad x\in\R .
  \label{eq:un-escalar}
\end{equation}
La razón para empezar aquí no es que el caso escalar sea útil por sí mismo, aunque lo es (el diodo túnel, una población con cosecha y el juego halcón--paloma de este capítulo son escalares), sino que en una dimensión \emph{todo se ve}. Las ideas que ocupan el resto del curso aparecen aquí sin el ruido de la geometría del plano: equilibrio, estabilidad local, linealización y sus límites, región de atracción, función potencial (un primer encuentro con las funciones de Lyapunov), conjunto invariante, bifurcación, existencia y unicidad, escape en tiempo finito. Cada una se volverá a encontrar, generalizada.

\notamargen{\textbf{Ruta de las clases.} Clase~1: \S\ref{sec:un-flujo}--\ref{sec:un-monotonia}. Clase~2: \S\ref{sec:un-poblaciones}--\ref{sec:un-escape}. Las dos últimas secciones motivan el \cref{cap:preliminares}.}

\begin{pausa}
Antes de leer, apueste: (1) ¿Puede oscilar un sistema $\dot x=f(x)$ con $x\in\R$? (2) Si $f(\eq)=0$ y $f'(\eq)=0$, ¿qué se concluye sobre $\eq$? (3) ¿Cuánto vale en $t=2$ la solución de $\dot x=x^2$, $x(0)=1$? (4) ¿Pueden dos soluciones distintas de $\dot x=f(x)$ salir del mismo punto?
\respuesta{(1) No, si $f$ es localmente Lipschitz: las soluciones son monótonas (\cref{prop:un-monotonia}). Sí, si el sistema es variante en el tiempo o si el estado vive en una circunferencia (\cref{tr:un-oscilar}). (2) Nada: puede ser estable, inestable o semiestable (\cref{tr:un-nohiperbolico}). (3) No existe: la solución escapó a infinito en $t=1$ (\cref{ej:un-escape}). (4) Sí, si $f$ no es Lipschitz, como en $\dot x=x^{1/3}$ (\cref{ej:un-raiz}); no, si lo es.}
\end{pausa}

%% ---------------------------------------------------------------------
\section{El flujo sobre la recta}
\label{sec:un-flujo}
%% ---------------------------------------------------------------------

La idea central, que viene de Poincaré y que \citet{Str15} explota de manera sistemática, es leer \eqref{eq:un-escalar} como un \clave{campo vectorial} sobre la recta: en cada punto $x$ la ecuación prescribe una velocidad $f(x)$. Imagine una partícula que se mueve por el eje real: donde $f(x)>0$ avanza a la derecha, donde $f(x)<0$ a la izquierda, y donde $f(x)=0$ se queda quieta. Basta graficar $f$ contra $x$ y marcar el sentido del movimiento sobre el eje para saber, sin resolver nada, qué hace cada solución. Esa gráfica es la \clave{línea de fase} del sistema.
\notamargen{\textbf{Convención} (la misma en las notas, en \cmd{lineaFase.m} y en la página web): círculo lleno, equilibrio estable; vacío, inestable; gris, semiestable; triángulos, sentido del flujo.}

\begin{ejemplo}[label=ej:un-seno]{El campo $\dot x=\sin x$}
Resuelva $\dot x=\sin x$ con $x(0)=\pi/4$ y diga qué pasa cuando $t\to\infty$.
\solucion Separando variables,
\[
  \int_{x_0}^{x}\frac{\dif s}{\sin s}=t\quad\Longrightarrow\quad\tan\frac{x(t)}2=\tan\frac{x_0}2\;e^{t},
\]
y para $x_0\in(-\pi,\pi)$, $x(t)=2\arctan\bigl(\tan(x_0/2)\,e^t\bigr)$. Con $x_0=\pi/4$, $\tan(\pi/8)=\sqrt2-1$ y
\[
  x(t)=2\arctan\bigl((\sqrt2-1)\,e^t\bigr)\xrightarrow[t\to\infty]{}2\cdot\frac\pi2=\pi ;
\]
por ejemplo, $x(1)=1.689$. La fórmula es exacta, pero no se lee de un vistazo. La línea de fase (\cref{fig:un-seno}) dice lo mismo sin cuentas: entre $0$ y $\pi$ el campo es positivo, así que la solución avanza a la derecha y se acerca a $\pi$, donde el flujo se detiene. Los equilibrios son $\eq=k\pi$, $k\in\Z$: los múltiplos pares de $\pi$ repelen y los impares atraen.
\end{ejemplo}

\begin{figure}[ht]
\centering
\includegraphics[width=\linewidth]{cap02/f_seno}
\caption{Izquierda: línea de fase de $\dot x=\sin x$. Derecha: quince soluciones; los círculos son la solución exacta desde $x_0=\pi/4$ (el error de \cmd{ode45} es menor que $10^{-9}$). Ninguna solución cruza las rectas $x=k\pi$.}
\label{fig:un-seno}
\end{figure}

\codigorepo{cap02_unidimensionales/C2_Ej1_Seno.m}{exacta}{Solución exacta contra \cmd{ode45}}

\video{V2_1_linea_fase}{La partícula sobre la recta}
La \cref{fig:un-seno} muestra un rasgo que conviene subrayar: \alerta{las soluciones nunca cruzan un equilibrio}. No es una casualidad numérica: es consecuencia de la unicidad de las soluciones, y es la clave de la \cref{sec:un-monotonia}.

\begin{trampa}[label=tr:un-infinito]{Lo que crece, crece sin límite}
\plausible{En $\dot x=\sin x$ con $x_0=\pi/4$, como $\dot x>0$, la solución crece para siempre: $x(t)\to\infty$.}
Así lo afirmaban las notas de 2024. Es cierto que $x(t)$ crece para siempre, pero crecer no es tender a infinito: la solución queda atrapada entre los equilibrios $0$ y $\pi$ y tiende a $\pi$. Además, $\abs{\sin x-\sin y}\le\abs{x-y}$, de modo que $\sin x$ es globalmente Lipschitz y las soluciones existen para todo $t$ sin escapar (\cref{sec:un-escape}).
\correcto{Una función creciente y acotada tiene límite. Antes de concluir que una solución diverge, busque el primer equilibrio en la dirección del flujo.}
\end{trampa}

%% ---------------------------------------------------------------------
\section{Estabilidad y linealización}
\label{sec:un-linealizacion}
%% ---------------------------------------------------------------------

La línea de fase sugiere tres clases de equilibrios aislados. Las definiciones precisas ($\varepsilon$--$\delta$) son del \cref{cap:lyapunov}; por ahora bastan las geométricas.

\begin{definicion}[label=def:un-clasificacion]{Clasificación en la recta}
Sea $\eq$ un equilibrio aislado de \eqref{eq:un-escalar}. Se dice que $\eq$ es \clave{estable} (un \emph{atractor}) si $f>0$ inmediatamente a la izquierda de $\eq$ y $f<0$ inmediatamente a la derecha; \clave{inestable} (un \emph{repulsor}) si los signos son los opuestos; y \clave{semiestable} si $f$ tiene el mismo signo a ambos lados.
\end{definicion}

Para decidir el tipo sin graficar se linealiza. Si $f\in\Cont1$ y $x=\eq+\eta$ con $\eta$ pequeño,
\[
  \dot\eta=f(\eq+\eta)=\underbrace{f(\eq)}_{=0}+f'(\eq)\,\eta+o(\eta) .
\]
Si $f'(\eq)\neq0$ el término lineal domina cerca de $\eq$, y la perturbación evoluciona como $\eta(t)\approx\eta(0)\,e^{f'(\eq)t}$.

\begin{proposicion}[label=prop:un-derivada]{Criterio de la derivada}
Sea $f\in\Cont1$ con $f(\eq)=0$. Si $f'(\eq)<0$, el equilibrio $\eq$ es estable; si $f'(\eq)>0$, es inestable. En ambos casos $\tau=1/\abs{f'(\eq)}$ es la escala de tiempo de la dinámica cerca de $\eq$.
\end{proposicion}

\begin{proof}
Si $f'(\eq)<0$, por la continuidad de $f'$ existe $r>0$ con $f'<0$ en $(\eq-r,\eq+r)$. Allí $f$ es estrictamente decreciente y, como $f(\eq)=0$, $f>0$ a la izquierda y $f<0$ a la derecha. El caso $f'(\eq)>0$ es simétrico.
\end{proof}

Un equilibrio con $f'(\eq)\neq0$ se llama \clave{hiperbólico}. Para $\sin x$, $f'(k\pi)=\cos k\pi=(-1)^k$: confirma la clasificación del \cref{ej:un-seno}, con $\tau=1$.

\begin{trampa}[label=tr:un-nohiperbolico]{$f'(\eq)=0$ no decide nada}
\plausible{Si $f'(\eq)=0$, la linealización $\dot\eta=0$ dice que las perturbaciones ni crecen ni decrecen: $\eq$ es marginalmente estable.}
La \cref{prop:un-derivada} no dice nada cuando $f'(\eq)=0$, y puede pasar cualquier cosa. Los sistemas
\[
  \dot x=-x^3,\qquad \dot x=x^3,\qquad \dot x=x^2
\]
tienen $f(0)=f'(0)=0$, la misma linealización, y el origen es estable, inestable y semiestable, respectivamente (\cref{fig:un-nohiperbolicos}). Hay más. En el caso estable la convergencia \emph{no es exponencial}: $x(t)=x_0/\sqrt{1+2x_0^2t}$ decae como $t^{-1/2}$ (con $x_0=1$, $x(10)=1/\sqrt{21}=0.218$). En los otros dos, las soluciones no solo se alejan del origen: \emph{escapan a infinito en tiempo finito} (\cref{sec:un-escape}).
\correcto{Un equilibrio no hiperbólico se clasifica con términos de orden superior o con un argumento global, nunca con la linealización. Este caso límite es el corazón de la teoría de bifurcaciones (capítulo~6).}
\end{trampa}

\begin{figure}[ht]
\centering
\includegraphics[width=\linewidth]{cap02/f_no_hiperbolicos}
\caption{Tres equilibrios no hiperbólicos con la misma linealización. Abajo, las soluciones se detienen al llegar a $\abs x=3$. La línea punteada es $1.5/\sqrt{1+4.5\,t}$.}
\label{fig:un-nohiperbolicos}
\end{figure}

\codigorepo{cap02_unidimensionales/C2_Ej2_NoHiperbolicos.m}{tres}{Tres sistemas con $f'(0)=0$}

%% ---------------------------------------------------------------------
\section{Por qué no hay oscilaciones en una dimensión}
\label{sec:un-monotonia}
%% ---------------------------------------------------------------------

La línea de fase deja ver algo más fuerte que la clasificación local: sobre la recta, una solución solo puede avanzar en un sentido.

\begin{proposicion}[label=prop:un-monotonia]{Monotonía de las soluciones}
Sea $f$ localmente Lipschitz (por ejemplo, $f\in\Cont1$) y sea $x(t)$ la solución de \eqref{eq:un-escalar} en su intervalo máximo de existencia $[0,T)$.
\begin{enumerate}
  \item Si $f(x_0)=0$, entonces $x(t)\equiv x_0$.
  \item Si $f(x_0)>0$ (resp. $<0$), entonces $x(t)$ es estrictamente creciente (resp. decreciente) y $f(x(t))\neq0$ para todo $t\in[0,T)$.
  \item Si además $x(t)$ es acotada, entonces $T=\infty$ y $x(t)$ converge a un equilibrio cuando $t\to\infty$.
\end{enumerate}
\end{proposicion}

\begin{proof}
(1) La constante $x_0$ es una solución y, por la unicidad (\cref{thm:pm-local}), es la única. (2) Si existiera $t_1$ con $f(x(t_1))=0$, la constante $x(t_1)$ y la solución $x(t)$ resolverían el mismo problema de valor inicial en $t_1$, lo que contradice la unicidad. Luego $f(x(t))$ no se anula y, por continuidad, conserva el signo de $f(x_0)$; como $\dot x=f(x)$, sigue la monotonía. (3) Que $T=\infty$ lo da el \cref{thm:pm-compacto}. Una función monótona y acotada tiene límite $\ell$. Si $f(\ell)\neq0$, por ejemplo $f(\ell)>0$, entonces $\dot x\ge f(\ell)/2$ para $t$ grande y $x$ crecería sin cota, lo que es absurdo. Luego $f(\ell)=0$.
\end{proof}

La consecuencia es inmediata: \alerta{un sistema escalar invariante en el tiempo, con $f$ localmente Lipschitz, no tiene soluciones periódicas no constantes}. Las oscilaciones exigen al menos dos dimensiones (capítulo~4) o un estado que viva en una circunferencia en lugar de una recta (\cref{cap:orbitas}).

\begin{trampa}[label=tr:un-oscilar]{En una dimensión nada oscila}
\plausible{Un sistema con una sola variable de estado no puede oscilar.}
La conclusión anterior depende de tres hipótesis, y falla si se quita cualquiera:
\begin{itemize}
  \item \emph{Invariancia en el tiempo.} $\dot x=-x+\sin t$ es escalar y sus soluciones tienden a la oscilación $\tfrac12(\sin t-\cos t)$.
  \item \emph{Estado en la recta.} $\dot\theta=\omega-\sin\theta$ con $\theta$ un ángulo tiene soluciones periódicas si $\omega>1$: la partícula da vueltas a la circunferencia.
  \item \emph{Unicidad.} Sin la condición de Lipschitz una solución puede detenerse en un equilibrio en tiempo finito, como $\dot x=-x^{1/3}$ (\cref{ej:un-raiz}).
\end{itemize}
\correcto{``No hay oscilaciones'' vale para $\dot x=f(x)$, $x\in\R$, con $f$ localmente Lipschitz. Las tres hipótesis cuentan.}
\end{trampa}

\subsection{La bolita que baja: el potencial}
\label{sec:un-potencial}

Hay otra manera de ver la \cref{prop:un-monotonia}, que anticipa el método directo de Lyapunov. Defina el \clave{potencial}
\begin{equation}
  V(x)\defeq-\int_0^xf(s)\,\dif s,\qquad\text{de modo que}\qquad\dot x=f(x)=-V'(x) .
  \label{eq:un-potencial}
\end{equation}
A lo largo de una solución,
\[
  \frac{\dif}{\dif t}V\bigl(x(t)\bigr)=V'(x)\,\dot x=-f(x)^2\le0 ,
\]
con igualdad solo en los equilibrios. La partícula siempre ``baja'' por la gráfica de $V$, como una bolita en un medio muy viscoso: los mínimos locales de $V$ son los equilibrios estables y los máximos, los inestables. Para $\dot x=\sin x$, $V(x)=\cos x-1$ (\cref{fig:un-potencial}). En el capítulo~7 el papel de $V$ lo hará una \clave{función de Lyapunov}: una función que decrece a lo largo de las soluciones aunque el sistema no sea un gradiente.

\begin{figure}[ht]
\centering
\includegraphics[width=\linewidth]{cap02/f_potencial}
\caption{Izquierda: potencial $V(x)=\cos x-1$ de $\dot x=\sin x$; los mínimos ($x=\pm\pi$) son los equilibrios estables. Derecha: $V(x(t))$ a lo largo de ocho soluciones; nunca crece.}
\label{fig:un-potencial}
\end{figure}

\codigorepo{cap02_unidimensionales/C2_Ej3_Potencial.m}{potencial}{El potencial a lo largo de las soluciones}

\begin{pausa}
¿Todo sistema escalar $\dot x=f(x)$ con $f$ continua tiene un potencial? ¿Y todo sistema en $\R^2$?
\respuesta{En $\R$ sí: $V=-\int f$ siempre existe. En $\R^2$ no: $\dot x=-\nabla V$ exige que el campo sea un gradiente, y un campo con rotación (como $\dot x_1=-x_2$, $\dot x_2=x_1$) no lo es. Por eso en dimensión mayor hay que \emph{buscar} una función que decrezca: es el problema de construir funciones de Lyapunov (capítulos~7 y~8).}
\end{pausa}

%% ---------------------------------------------------------------------
\section{Poblaciones: la logística y la cosecha}
\label{sec:un-poblaciones}
%% ---------------------------------------------------------------------

La ecuación logística del \cref{cap:introduccion}, reescalada con $x=P/K$ y $\tau=Rt$ (\cref{prob:in-reescalar}), es $\dot x=x(1-x)$: equilibrios en $0$ (inestable, $f'(0)=1$) y $1$ (estable, $f'(1)=-1$). Toda población positiva crece monótonamente hasta la capacidad del medio. Ahora se le quita, cada unidad de tiempo, una cantidad fija $h\ge0$ (pesca, caza, tala), medida en las mismas unidades reescaladas:
\begin{equation}
  \dot x=f_h(x)\defeq x(1-x)-h .
  \label{eq:un-cosecha}
\end{equation}

\begin{ejemplo}[label=ej:un-cosecha]{Cosecha sostenible y colapso}
Analice \eqref{eq:un-cosecha} para $h=0.2$ y $h=0.3$. ¿Cuál es la cosecha máxima que la población resiste?
\solucion Los equilibrios son las raíces de $x^2-x+h=0$, $x^*_\pm=\tfrac12\bigl(1\pm\sqrt{1-4h}\bigr)$, reales si $h\le1/4$. Con $h=0.2$: $x^*_-=0.276$ y $x^*_+=0.724$, con $f_h'(x^*_\pm)=1-2x^*_\pm=\mp0.447$: el de arriba es estable y el de abajo inestable. El de abajo es un \emph{umbral}: si la población cae por debajo de $0.276$ (por una mala temporada, por ejemplo), se extingue aunque la cosecha no cambie. Con $h=0.3>1/4$ no hay equilibrios: $f_h<0$ para todo $x$ y la población llega a cero en tiempo finito desde cualquier condición inicial. Desde $x_0=0.5$, con $a=\sqrt{h-1/4}$, la solución es $x(t)=\tfrac12-a\tan(at)$ y llega a $0$ en $t=\arctan(1/(2a))/a=5.14$ (\cref{fig:un-cosecha}). La cosecha crítica es $h_c=1/4$, el máximo de $x(1-x)$: en unidades originales, $H_c=RK/4$.
\end{ejemplo}

\begin{figure}[ht]
\centering
\includegraphics[width=\linewidth]{cap02/f_cosecha}
\caption{Izquierda: $f_h(x)=x(1-x)-h$ para cuatro cosechas; los dos equilibrios se acercan, chocan en $h=1/4$ (gris) y desaparecen. Derecha: soluciones con $h=0.2$ (azul; las punteadas son los equilibrios) y con $h=0.3$ desde $x_0=0.5$ (coral).}
\label{fig:un-cosecha}
\end{figure}

\codigorepo{cap02_unidimensionales/C2_Ej4_Logistica.m}{cosecha}{Equilibrios de la logística con cosecha}

Al subir $h$, los dos equilibrios se acercan, chocan en $h=1/4$ y desaparecen: es la \clave{bifurcación silla-nodo} del \cref{cap:introduccion} ($\dot x=r+x^2$), que se estudia en el capítulo~6. Y deja una lección de manejo de recursos: el punto de máxima cosecha sostenible está \emph{justo en el borde del colapso}. Operar allí es operar sin margen.

\begin{trampa}[label=tr:un-sostenible]{Cosechar al máximo sostenible es lo óptimo}
\plausible{La cosecha óptima es $h=1/4$: es la mayor para la que existe un equilibrio.}
En $h=1/4$ el único equilibrio, $x=1/2$, es semiestable ($f_h\le0$ a ambos lados): cualquier perturbación hacia abajo lleva a la extinción. Además, cualquier error en la estimación de $R$ o de $K$ pone a la población del lado $h>h_c$, donde el colapso ocurre en tiempo finito.
\correcto{El análisis de un equilibrio incluye su \emph{robustez}: qué tan cerca está de desaparecer y qué tan grande es su región de atracción. En $h=0.2$ esa región es $(0.276,\infty)$; en $h=1/4$ es solo $[1/2,\infty)$.}
\end{trampa}

%% ---------------------------------------------------------------------
\section{Un circuito con memoria: el diodo túnel}
\label{sec:un-diodo}
%% ---------------------------------------------------------------------

El diodo túnel tiene una característica corriente--voltaje $i=h(v)$ con una región de \clave{resistencia diferencial negativa}: en un intervalo de voltajes, la corriente \emph{baja} cuando el voltaje sube. Conectado en paralelo con un condensador $C$ y alimentado por una fuente $E$ a través de una resistencia $R$ (\cref{fig:un-circuito}), la ley de corrientes de Kirchhoff en el nodo da
\begin{equation}
  C\,\dot v=\frac{E-v}{R}-h(v) .
  \label{eq:un-diodo}
\end{equation}
\notamargen{Es el circuito de la sección~1.2.2 de \citet{Kha02} sin la inductancia; con ella el modelo es de segundo orden (\cref{prob:in-diodo} y \cref{sec:bi-ejemplos}). Unidades: $v$ en V, $h$ en mA, $R$ en k$\Omega$, $C$ en pF y $t$ en ns, porque $1\,\text{mA}/1\,\text{pF}=1\,\text{V/ns}$.}
Se usa la característica de \citet{Kha02},
\[
  h(v)=17.76v-103.79v^2+229.62v^3-226.31v^4+83.72v^5 ,
\]
con $E=1.2$\,V, $R=1.5$\,k$\Omega$ y $C=2$\,pF.

\begin{figure}[ht]
\centering
\includegraphics[width=0.8\linewidth]{cap02/t_diodo}
\caption{Circuito con diodo túnel. El estado es el voltaje $v$ del condensador.}
\label{fig:un-circuito}
\end{figure}

Los equilibrios cumplen $h(v)=(E-v)/R$: son las intersecciones de la característica del diodo con la \clave{recta de carga} (\cref{fig:un-diodo}). Hay tres. Derivando el lado derecho de \eqref{eq:un-diodo},
\[
  f'(\eq)=-\frac1C\Bigl(\frac1R+h'(\eq)\Bigr),
\]
de modo que un equilibrio es inestable exactamente cuando la resistencia diferencial negativa del diodo vence la conductancia de la carga, $h'(\eq)<-1/R$. Eso solo ocurre en el del medio (\cref{tab:un-diodo}).

\begin{table}[ht]
\centering
\begin{tabular}{@{}ccccl@{}}
\toprule
$v^*$ [V] & $i^*$ [mA] & $f'(v^*)$ [ns$^{-1}$] & $\tau$ [ns] & tipo\\
\midrule
$0.0626$ & $0.758$ & $-3.956$ & $0.253$ & estable\\
$0.2854$ & $0.610$ & $+1.487$ & $0.673$ & inestable\\
$0.8844$ & $0.210$ & $-1.770$ & $0.565$ & estable\\
\bottomrule
\end{tabular}
\caption{Equilibrios del circuito con diodo túnel ($E=1.2$\,V). Coinciden con los de \citet{Kha02}.}
\label{tab:un-diodo}
\end{table}

\begin{figure}[ht]
\centering
\includegraphics[width=\linewidth]{cap02/f_diodo}
\caption{Izquierda: característica del diodo y recta de carga; sus intersecciones son los equilibrios. Derecha: línea de fase de \eqref{eq:un-diodo}.}
\label{fig:un-diodo}
\end{figure}

\codigorepo{cap02_unidimensionales/C2_Ej5_DiodoTunel.m}{equilibrios}{Equilibrios del diodo túnel y su clasificación}

El circuito es \clave{biestable} (\cref{fig:un-diodo-tiempo}): toda condición inicial por debajo de $0.2854$\,V termina en el estado ``bajo'' y toda condición por encima, en el ``alto''. El equilibrio inestable es la frontera entre las dos \clave{regiones de atracción}: $(-\infty,\,0.2854)$ para el bajo y $(0.2854,\,\infty)$ para el alto. Es una celda de memoria: el circuito ``recuerda'' de qué lado del umbral se le dejó. Ningún circuito lineal puede hacerlo, porque no puede tener dos equilibrios aislados.

\begin{figure}[ht]
\centering
\includegraphics[width=0.7\linewidth]{cap02/f_diodo_tiempo}
\caption{Respuestas $v(t)$ para $v(0)=0,0.1,\dots,1$\,V. En menos de 3\,ns cada una llega al equilibrio estable de su lado del umbral.}
\label{fig:un-diodo-tiempo}
\end{figure}

\subsection{Histéresis}

¿Qué pasa si el voltaje $E$ de la fuente cambia lentamente? Los equilibrios son los puntos de la curva $E=v+R\,h(v)$, que tiene un máximo local en $v=0.149$ ($E=1.643$\,V) y un mínimo local en $v=0.513$ ($E=0.659$\,V): \alerta{hay tres equilibrios para $0.659<E<1.643$\,V} y uno solo fuera de ese intervalo. Si $E$ sube desde $0$, el voltaje sigue la rama baja; cuando la rama baja desaparece, en $E\approx1.64$\,V, el estado no tiene a dónde ir más que a la rama alta, y salta. Si después $E$ baja, el estado se queda en la rama alta hasta que esta desaparece, en $E\approx0.66$\,V, y salta abajo (\cref{fig:un-histeresis}). El estado depende de la historia de $E$, no solo de su valor actual: eso es \clave{histéresis}. Cada salto ocurre cuando un equilibrio estable choca con el inestable y ambos desaparecen: dos bifurcaciones silla-nodo.


\codigorepo{cap02_unidimensionales/C2_Ej6_Histeresis.m}{intervalo}{¿Para qué $E$ hay tres equilibrios?}
\codigorepo{cap02_unidimensionales/C2_Ej6_Histeresis.m}{barrido}{Barrido lento de $E$: los saltos}

\begin{figure}[!htbp]
\centering
\includegraphics[width=\linewidth]{cap02/f_histeresis}
\caption{Izquierda: $E$ sube de $0$ a $1.8$\,V en 4000\,ns y vuelve a bajar; el estado (azul al subir, coral al bajar) recorre las ramas estables de la curva de equilibrios y salta en sus extremos. Derecha: el lado derecho de \eqref{eq:un-diodo} con uno, tres y un equilibrio.}
\label{fig:un-histeresis}
\end{figure}

\video{V2_2_histeresis}{El diodo túnel recuerda}
\begin{modelo}[DiodoTunel.slx]
\cmd{crearModelo\_DiodoTunel.m} arma \eqref{eq:un-diodo} con un integrador, un bloque \emph{Polynomial} para $h(v)$ y una fuente que sube y baja; simula 8000\,ns y compara con \cmd{ode45}.
\end{modelo}

\begin{trampa}[label=tr:un-histeresis]{El salto ocurre exactamente en el valor estático}
\plausible{En el barrido, el voltaje salta a la rama alta justo cuando $E=1.643$\,V.}
En la simulación salta en $E=1.652$\,V, y al bajar en $0.650$\,V en lugar de $0.659$. No es un error numérico: cerca de una silla-nodo la dinámica es muy lenta (el campo es casi tangente al eje, $f'\approx0$), y el estado se demora en abandonar la rama que acaba de desaparecer. Mientras más rápido el barrido, mayor el retraso.
\correcto{El análisis de equilibrios (cuasiestático) da los valores a los que tiende el salto cuando el barrido es infinitamente lento. Cerca de una bifurcación, la escala de tiempo $1/\abs{f'(\eq)}$ crece sin límite y la separación de escalas deja de valer.}
\end{trampa}

%% ---------------------------------------------------------------------
\section{Halcones y palomas: una dinámica evolutiva}
\label{sec:un-halcon}
%% ---------------------------------------------------------------------

El último ejemplo viene de la teoría de juegos evolutiva y muestra que la línea de fase no depende del origen físico del modelo.

\begin{historia}[La lógica del conflicto animal]
En 1973, John Maynard Smith y George Price se preguntaron por qué los conflictos entre animales de la misma especie rara vez escalan hasta el daño grave, aunque los contendientes tienen armas eficaces (cuernos, colmillos) \citep{MSPr73}. Su respuesta introdujo la noción de \emph{estrategia evolutivamente estable}. En 1978, Peter Taylor y Leo Jonker propusieron una dinámica que describe cómo cambia la composición de una población que usa varias estrategias \citep{TaJo78}; el nombre \emph{dinámica del replicador} lo acuñaron Schuster y Sigmund en 1983 \citep{ScSi83}. La versión de dos estrategias con el nombre halcón--paloma se popularizó con el libro de Maynard Smith \citep{MS82}.
\end{historia}

\paragraph{La dinámica del replicador.} Una población usa $N$ estrategias. Su estado es el vector de proporciones $p=(p_1,\dots,p_N)\T$, que vive en el \clave{símplex}
\[
  \simplex=\Bigl\{p\in\R^N:\ p_i\ge0,\ \textstyle\sum_{i=1}^Np_i=1\Bigr\}.
\]
En un juego simétrico con \clave{matriz de pagos} $\Pi\in\R^{N\times N}$, el pago esperado de la estrategia $i$ es $F_i(p)=(\Pi p)_i$ y el pago medio es $\bar F(p)=p\T\Pi p$. Una estrategia crece en proporción a cuánto supera su pago al promedio:
\begin{equation}
  \dot p_i=p_i\bigl(F_i(p)-\bar F(p)\bigr),\qquad i=1,\dots,N .
  \label{eq:un-replicador}
\end{equation}
\notamargen{Se usa $F$ para los pagos, como en \citet{San10}. Las notas de 2024 escribían $\omega(\bar\pi_i,p)$ y $J_1$; se cambió para no chocar con la frecuencia $\omega$ ni con los conjuntos $\omega$-límite.}

\paragraph{El juego.} Dos individuos se disputan un recurso de valor $v>0$. Un \emph{halcón} pelea hasta ganar o salir herido (costo $c>0$); una \emph{paloma} amenaza, pero se retira si el otro escala. Dos halcones ganan cada uno la mitad de las veces: pago esperado $(v-c)/2$. Un halcón contra una paloma se lleva $v$ y la paloma $0$. Dos palomas se reparten el recurso: $v/2$. La matriz de pagos del que elige la fila es
\begin{equation}
  \Pi_{\mathrm o}=\begin{bmatrix}\dfrac{v-c}2 & v\\[6pt] 0 & \dfrac v2\end{bmatrix}.
  \label{eq:un-hd}
\end{equation}

\begin{lema}[label=lem:un-columnas]{Invarianza por columnas}
Si a todos los elementos de una columna $j$ de $\Pi$ se les suma la misma constante $k$, la dinámica \eqref{eq:un-replicador} no cambia en $\simplex$.
\end{lema}

\begin{proof}
Con $\Pi'=\Pi+k\,\mathbf1e_j\T$, $F_i'(p)=F_i(p)+kp_j$ para todo $i$, y $\bar F'(p)=\bar F(p)+kp_j\sum_ip_i=\bar F(p)+kp_j$ porque $p\in\simplex$. La diferencia $F_i'-\bar F'$ no cambia.
\end{proof}

Restando $v$ a la segunda columna de $\Pi_{\mathrm o}$ queda $\Pi=\diag\bigl((v-c)/2,\,-v/2\bigr)$, que genera la misma dinámica. Con $p_2=1-p_1$, \eqref{eq:un-replicador} se reduce a un sistema escalar:
\begin{equation}
  \dot p_1=f(p_1)\defeq\tfrac12\,p_1(1-p_1)(v-c\,p_1),\qquad p_1\in[0,1] .
  \label{eq:un-hd-escalar}
\end{equation}

\begin{ejemplo}[label=ej:un-hd]{Halcones y palomas en la línea de fase}
Encuentre y clasifique los equilibrios de \eqref{eq:un-hd-escalar} según el signo de $v-c$.
\solucion Los equilibrios son $0$, $1$ y $v/c$. Derivando,
\[
  f'(p_1)=\tfrac12\bigl[(1-2p_1)(v-cp_1)-c\,p_1(1-p_1)\bigr],
\]
de donde $f'(0)=v/2>0$, $f'(1)=(c-v)/2$ y $f'(v/c)=-\tfrac v2\bigl(1-\tfrac vc\bigr)$.
\begin{itemize}
  \item Si $v<c$ (pelear cuesta más de lo que se gana): los tres están en $[0,1]$; $0$ y $1$ son inestables y $v/c$ es estable. La población converge a una mezcla con una fracción $v/c$ de halcones desde cualquier estado interior (\cref{fig:un-halcon}). Con $v=1$ y $c=3$, a $1/3$, con $f'(1/3)=-1/3$. Esa es la respuesta de Maynard Smith y Price: no todos los encuentros escalan porque, en equilibrio, solo una parte de la población usa la estrategia agresiva.
  \item Si $v>c$: $v/c>1$ no está en el espacio de estado; $f>0$ en $(0,1)$ y la población converge a $p_1=1$: todos halcones.
\end{itemize}
\end{ejemplo}

\begin{figure}[ht]
\centering
\includegraphics[width=\linewidth]{cap02/f_halcon}
\caption{Línea de fase de \eqref{eq:un-hd-escalar} con $v<c$ y con $v>c$, y fracción de halcones con $v=1$, $c=3$ desde siete condiciones iniciales.}
\label{fig:un-halcon}
\end{figure}

\codigorepo{cap02_unidimensionales/C2_Ej7_HalconPaloma.m}{replicador}{Replicador: matriz original, matriz reducida y forma escalar}

Dos hechos de este ejemplo reaparecerán. Primero, el intervalo $[0,1]$ es \clave{invariante}: como $0$ y $1$ son equilibrios, la \cref{prop:un-monotonia} impide que una solución que empieza en $[0,1]$ lo abandone. Es el primer ejemplo de un conjunto invariante, la idea central del principio de LaSalle (capítulo~8). Segundo, el sistema escalar se obtuvo restringiendo un sistema de dimensión 2 a la recta invariante $p_1+p_2=1$; en el \cref{cap:pasividad} la misma dinámica del replicador reaparece repartiendo agua entre tanques.

\begin{trampa}[label=tr:un-dominio]{Los equilibrios son las raíces de $f$}
\plausible{Los equilibrios del juego halcón--paloma son siempre $0$, $1$ y $v/c$.}
Las tres raíces de $f$ existen siempre, pero $v/c$ es un estado de la población solo si $v/c\le1$. Si $v>c$, la raíz $v/c$ es un número sin significado: una proporción mayor que uno.
\correcto{Una raíz de $f$ es un equilibrio \emph{del modelo} solo si está en el conjunto donde el modelo tiene sentido. Antes de clasificar equilibrios hay que verificar que estén en el dominio y que el dominio sea invariante.}
\end{trampa}

%% ---------------------------------------------------------------------
\section{¿Es única la solución?}
\label{sec:un-unicidad}
%% ---------------------------------------------------------------------

Todo lo anterior supuso, sin decirlo, que por cada condición inicial pasa exactamente una solución; sin eso, la línea de fase pierde el sentido. Una solución de \eqref{eq:un-escalar} en $[0,T)$ es una función continuamente diferenciable que cumple la ecuación; equivalentemente, una función continua que cumple la \clave{ecuación integral}
\begin{equation}
  x(t)=x_0+\int_0^tf\bigl(x(s)\bigr)\,\dif s .
  \label{eq:un-integral}
\end{equation}
\notamargen{Atención al integrando: es $f(x(s))$, no $f(s)$. La solución aparece a los dos lados: \eqref{eq:un-integral} es una ecuación, no una fórmula.}
Si $f$ es continua, siempre existe al menos una solución local (teorema de Peano, \citealp{Pea1890}). La unicidad pide más.

\begin{ejemplo}[label=ej:un-raiz]{Infinitas soluciones desde el mismo punto}
Considere $\dot x=x^{1/3}$, $x(0)=0$, con $x^{1/3}$ la raíz cúbica real. ¿Cuántas soluciones tiene?
\solucion $x(t)\equiv0$ es una. Separando variables con $x>0$, $\tfrac32x^{2/3}=t-c$, y para cada $c\ge0$
\[
  x_c(t)=\begin{cases}0, & 0\le t\le c,\\[2pt] \bigl(\tfrac23(t-c)\bigr)^{3/2}, & t>c,\end{cases}
\]
también lo es: para $t>c$, $\dot x_c=\bigl(\tfrac23(t-c)\bigr)^{1/2}=x_c^{1/3}$, y las dos ramas empalman con derivada cero. Lo mismo con $-x_c$. Por el origen pasan infinitas soluciones: la partícula puede quedarse quieta un tiempo arbitrario y luego ``arrancar'' sola (\cref{fig:un-nounicidad}). El sistema emparentado $\dot x=-x^{1/3}$ hace lo contrario: desde $x_0>0$ llega a $0$ en el tiempo finito $T=\tfrac32x_0^{2/3}$ y se queda ahí, algo imposible para un $f$ localmente Lipschitz (si se corre el tiempo hacia atrás desde el equilibrio, saldrían dos soluciones del mismo punto).
\end{ejemplo}

\begin{figure}[ht]
\centering
\includegraphics[width=\linewidth]{cap02/f_no_unicidad}
\caption{Izquierda: algunas de las infinitas soluciones de $\dot x=x^{1/3}$, $x(0)=0$; \cmd{ode45} devuelve la trivial. Derecha: $\dot x=-x^{1/3}$ llega a cero en tiempo finito.}
\label{fig:un-nounicidad}
\end{figure}

¿Qué falla? La pendiente de $x^{1/3}$ es $\tfrac13x^{-2/3}$, que no está acotada cerca del origen. Lo que se necesita es controlar cuánto puede cambiar $f$ entre dos puntos cercanos: $f$ es \clave{Lipschitz} en un conjunto $W$ si existe $L$ con
\[
  \abs{f(x)-f(y)}\le L\,\abs{x-y}\qquad\text{para todo }x,y\in W ,
\]
y \clave{localmente Lipschitz} si cada punto tiene una vecindad donde lo es. Toda $f\in\Cont1$ es localmente Lipschitz (por el teorema del valor medio, con $L=\max\abs{f'}$ en un intervalo cerrado); $\abs x$ es Lipschitz sin ser derivable; $x^{1/3}$ es continua y no es Lipschitz en ninguna vecindad del origen. El resultado que se usa en todo el curso es:
\begin{center}
\fbox{\parbox{0.9\linewidth}{\centering si $f$ es localmente Lipschitz, por cada $x_0$ pasa una única solución, definida en un intervalo $[0,\delta)$.}}
\end{center}
Su enunciado preciso y su demostración (con el principio de contracción) son el \cref{thm:pm-local}. La condición lleva el nombre de Rudolf Lipschitz, que la usó en 1876 para probar la unicidad \citep{Lip1876}.
\notamargen{En Francia el resultado se conoce como teorema de Cauchy--Lipschitz; en buena parte de la literatura, como teorema de Picard--Lindelöf, por la prueba con aproximaciones sucesivas \citep{Pic1890,Lin1894}.}

\begin{trampa}[label=tr:un-raiz]{El simulador me dio una sola solución}
\plausible{Integré $\dot x=x^{1/3}$ desde $x_0=0$ con \cmd{ode45} y obtuve $x\equiv0$: la solución es única.}
Desde $x_0=10^{-12}$, que en cualquier experimento es cero, \cmd{ode45} sigue la rama $x_0$ y llega a $2\sqrt2=2.83$ en $t=3$. \emph{Un integrador siempre devuelve una sola trayectoria}: no puede detectar la falta de unicidad, que es una propiedad del problema y no del algoritmo. Y hay una trampa de implementación: en MATLAB, \cmd{(-8)\^{}(1/3)} devuelve el número complejo $1+1.7321\,\ii$ (la raíz principal), no $-2$. La raíz cúbica real es \cmd{nthroot(x,3)}.
\correcto{Las hipótesis de unicidad se verifican con lápiz y papel antes de simular. Si fallan, el modelo predice más de un futuro y hay que decidir cuál es el físicamente relevante.}
\end{trampa}

\codigorepo{cap02_unidimensionales/C2_Ej8_NoUnicidad.m}{raiz}{La raíz cúbica en MATLAB}
\codigorepo{cap02_unidimensionales/C2_Ej8_NoUnicidad.m}{tiempofinito}{Llegar a un equilibrio en tiempo finito}

%% ---------------------------------------------------------------------
\section{Existencia global y escape en tiempo finito}
\label{sec:un-escape}
%% ---------------------------------------------------------------------

La unicidad local garantiza una solución en un intervalo $[0,\delta)$, quizás muy corto. Que $f$ sea suave no basta para que exista para todo $t\ge0$.

\begin{ejemplo}[label=ej:un-escape]{Escape en tiempo finito}
Resuelva $\dot x=x^2$, $x(0)=x_0>0$.
\solucion $f(x)=x^2$ es $\Cont\infty$ y, por tanto, localmente Lipschitz. Separando variables,
\[
  x(t)=\frac{x_0}{1-x_0t},\qquad 0\le t<\frac1{x_0} .
\]
La solución tiende a $+\infty$ cuando $t\to1/x_0$: el \clave{tiempo de escape} es $\bar t=1/x_0$ y el intervalo máximo de existencia es $[0,1/x_0)$. La fórmula vuelve a tener sentido para $t>1/x_0$, pero esa rama no está conectada con la condición inicial. Con $x_0=1$, la respuesta a la tercera pregunta de la \cref{sec:un-porque} es que la solución \emph{no existe} en $t=2$: dejó de existir en $t=1$.
\end{ejemplo}

¿Por qué escapa $x^2$ y no $\sin x$? Por la rapidez con que crece $f$ para $\abs x$ grande: $x^2$ crece más que linealmente, la velocidad aumenta tan deprisa que la partícula recorre una distancia infinita en un tiempo finito. Para descartar el escape hay dos herramientas: si $f$ es \emph{globalmente} Lipschitz, las soluciones existen para todo $t$ (\cref{thm:pm-global}); y si se sabe de antemano que la solución se queda en un compacto, también (\cref{thm:pm-compacto}). La primera es fácil de verificar y muy restrictiva; la segunda es la que se usa en la práctica, y así se usó en la \cref{prop:un-monotonia} y se usará con las funciones de Lyapunov: si $V$ no crece, las soluciones se quedan en un conjunto de nivel compacto y existen para siempre.

\begin{ejemplo}[label=ej:un-x2menos1]{Estable cerca de $-1$, explosivo después de $+1$}
Analice $\dot x=x^2-1$.
\solucion Los equilibrios son $-1$ (estable, $f'(-1)=-2$) y $1$ (inestable, $f'(1)=2$). Las notas de 2024 decían que una perturbación grande de $-1$ ``puede llevar al equilibrio $1$''. La línea de fase dice algo más preciso y más grave: si $x_0>1$, la solución no va a ningún equilibrio, escapa. Resolviendo,
\[
  \frac{x(t)-1}{x(t)+1}=\frac{x_0-1}{x_0+1}\,e^{2t}\quad\Longrightarrow\quad\bar t=\frac12\ln\frac{x_0+1}{x_0-1}\qquad(x_0>1) .
\]
Desde $x_0=1.5$ escapa en $\bar t=\tfrac12\ln5=0.805$; desde $x_0=3$, en $\tfrac12\ln2=0.347$. La región de atracción de $-1$ es exactamente $(-\infty,1)$ (\cref{fig:un-escape}), y la solución que empieza en $1$ se queda en $1$, pero la más mínima perturbación hacia arriba escapa.
\end{ejemplo}

\begin{figure}[ht]
\centering
\includegraphics[width=\linewidth]{cap02/f_escape}
\caption{Izquierda: $\dot x=x^2$ con cuatro condiciones iniciales; las líneas punteadas marcan $\bar t=1/x_0$. Derecha: $\dot x=x^2-1$; las soluciones con $x_0>1$ (coral) escapan.}
\label{fig:un-escape}
\end{figure}

\begin{trampa}[label=tr:un-suave]{Si $f$ es suave, la solución existe para siempre}
\plausible{$f(x)=x^2$ es un polinomio, infinitamente derivable: sus soluciones existen para todo $t$.}
La suavidad de $f$, incluso si es polinomial, garantiza existencia y unicidad \emph{locales}, nada más. En un sistema no lineal, ``inestable'' puede significar escapar a infinito en tiempo finito, algo imposible para $\dot x=Ax$, cuyas soluciones $e^{At}x_0$ existen siempre. Y en la simulación el escape es engañoso: \cmd{ode45} reduce el paso al acercarse a la asíntota y termina con la advertencia \emph{Unable to meet integration tolerances}, que un usuario desprevenido atribuye a un problema numérico.
\correcto{Antes de simular, pregúntese si $f$ crece más que linealmente. Si es así, detenga la integración con un evento cuando $\abs x$ supere una cota (opción \cmd{'Events'} de \cmd{odeset}), y compare el instante del evento con el tiempo de escape.}
\end{trampa}

\codigorepo{cap02_unidimensionales/C2_Ej9_Escape.m}{eventos}{Escape en tiempo finito con detección por eventos}

%% ---------------------------------------------------------------------
\section{Para llevar}
\label{sec:un-cierre}
%% ---------------------------------------------------------------------

\begin{parallevar}
\begin{enumerate}
  \item \textbf{La gráfica de $f$ lo dice todo} en una dimensión: los ceros son los equilibrios y el signo de $f$ da el sentido del flujo.
  \item \textbf{Linealización:} $f'(\eq)<0$, estable; $f'(\eq)>0$, inestable; $\tau=1/\abs{f'(\eq)}$. Si $f'(\eq)=0$, la linealización no decide.
  \item \textbf{Con unicidad, las soluciones escalares son monótonas:} no hay oscilaciones, ninguna solución cruza un equilibrio y toda solución acotada converge a uno.
  \item \textbf{El potencial} $V=-\int f$ decrece a lo largo de las soluciones: el primer ejemplo de función de Lyapunov.
  \item \textbf{Varios equilibrios estables} dan biestabilidad, regiones de atracción separadas por un equilibrio inestable y, si un parámetro cambia, histéresis.
  \item \textbf{Cuando los equilibrios chocan} (cosecha crítica, extremos de la histéresis) el sistema se reorganiza: silla-nodo.
  \item \textbf{Continuidad da existencia; Lipschitz da unicidad; ninguna da existencia global.} Un simulador devuelve una sola trayectoria aunque haya infinitas, y ``falla'' justo cuando la solución escapa.
\end{enumerate}
\end{parallevar}

\begin{trampa}[title={\faBug\hspace{0.45em}Las trampas del capítulo, en una caja}]{}
\begin{itemize}
  \item Crecer no es tender a infinito (\cref{tr:un-infinito}); $f'(\eq)=0$ no decide nada (\cref{tr:un-nohiperbolico}).
  \item ``En una dimensión nada oscila'' necesita tres hipótesis (\cref{tr:un-oscilar}).
  \item La cosecha máxima sostenible está en el borde del colapso (\cref{tr:un-sostenible}).
  \item Cerca de una silla-nodo el salto se retrasa (\cref{tr:un-histeresis}).
  \item Una raíz de $f$ fuera del dominio no es un equilibrio (\cref{tr:un-dominio}).
  \item El simulador no detecta la no unicidad, y \cmd{(-8)\^{}(1/3)} no es $-2$ (\cref{tr:un-raiz}).
  \item Suave no implica existencia global (\cref{tr:un-suave}).
\end{itemize}
\end{trampa}

\begin{profundiza}[Correcciones respecto a las versiones anteriores (Handout2, el borrador de LibroSisNL2026 y SisUniDim2026)]
(i) En $\dot x=\sin x$ con $x_0=\pi/4$ se afirmaba que $x(t)\to\infty$; la solución tiende a $\pi$ (\cref{tr:un-infinito}). (ii) En $\dot x=x^2-1$ se decía que una perturbación grande de $-1$ ``podría tender a $1$''; si $x_0>1$ la solución escapa en tiempo finito (\cref{ej:un-x2menos1}). (iii) La ecuación integral se escribía con $f(\tau)$ en lugar de $f(x(\tau))$. (iv) En el juego halcón--paloma se cambió la notación de los pagos ($F$ en lugar de $\omega(\bar\pi_i,p)$ y $J_1$) y se probó la invarianza por columnas, que antes se usaba sin justificar (\cref{lem:un-columnas}). (v) El tiempo de escape de $\dot x=x^2$ se llamaba \emph{finite escape time} sin decir que la solución no existe después; se aclaró que la rama de la fórmula para $t>1/x_0$ no es parte de la solución. (vi) Se agregaron la logística con cosecha y la histéresis del diodo túnel, que conectan con la silla-nodo del capítulo~6; las bifurcaciones de las diapositivas de 2026 (silla-nodo, transcrítica, horquilla, láser) quedan para ese capítulo. (vii) El teorema de existencia y unicidad se enuncia y se demuestra en el \cref{cap:preliminares}; aquí solo se usa. (viii) Notación de \texttt{notacion.sty}: invariante en el tiempo (no autónomo), $\eq$, $\Cont1$. Todos los números se verificaron con \texttt{herramientas/cap02\_numeros.py} (42 de 42) y \cmd{verificarCap02.m} (37 de 37).
\end{profundiza}

%% ---------------------------------------------------------------------
\section{Problemas}
\label{sec:un-problemas}
%% ---------------------------------------------------------------------

Los problemas marcados con \etiquetaM{} requieren \matlab; los marcados con \etiquetaEstrella{} son de profundización. Cada uno tiene una trampa: búsquela antes de dar la respuesta.

\begin{problema}[silla-nodo]\label{prob:un-sillanodo}
Trace la línea de fase de $\dot x=r+x^2$ para $r<0$, $r=0$ y $r>0$. ¿Cuántos equilibrios hay en cada caso? ¿Qué le pasa a $f'$ en el equilibrio estable cuando $r\to0^-$? ¿Qué significa eso para la velocidad de convergencia?
\end{problema}

\begin{problema}[efecto Allee]\label{prob:un-allee}
Analice $\dot x=x(1-x)(x-a)$, $0<a<1$, como modelo de una población que se extingue si es demasiado pequeña. Clasifique los equilibrios y dé la región de atracción de cada equilibrio estable. Compare con la cosecha (\cref{ej:un-cosecha}): ¿qué tienen en común $a$ y $x^*_-$?
\end{problema}

\begin{problema}[convergencia en tiempo finito]\label{prob:un-finito}
Muestre que la solución de $\dot x=-x^{1/3}$, $x(0)=x_0>0$, llega a $0$ en $T=\tfrac32x_0^{2/3}$ y se queda ahí. ¿Por qué no contradice la \cref{prop:un-monotonia}? Pruebe que ningún sistema $\dot x=f(x)$ con $f$ localmente Lipschitz llega a un equilibrio en tiempo finito.
\end{problema}

\begin{problema}[clasificar sin dibujar]\label{prob:un-clasificar}
Clasifique los equilibrios de (a) $\dot x=x^3-x$; (b) $\dot x=e^{-x}\sin x$; (c) $\dot x=x^2(1-x)$; (d) $\dot x=1-x-e^{-2x}$. ¿En cuál falla la \cref{prop:un-derivada} y qué hace en su lugar?
\end{problema}

\begin{problema}[un juego de coordinación]\label{prob:un-coordinacion}
Repita el \cref{ej:un-hd} con $\Pi=\diag(a,b)$, $a,b>0$ (las dos estrategias ganan cuando coinciden). Muestre que el equilibrio interior es inestable y dé la región de atracción de cada estrategia pura. Compare con el diodo túnel.
\end{problema}

\begin{problema}[histéresis]\label{prob:un-histeresis}
\etiquetaM{} (a) Con \cmd{C2\_Ej6\_Histeresis}, repita el barrido de $E$ diez veces más rápido. ¿Cuánto se retrasan los saltos? (b) ¿Para qué valores de $R$ desaparece la histéresis? (Pista: la curva $E=v+Rh(v)$ debe ser monótona; la condición es $1+Rh'(v)>0$ para todo $v$.)
\end{problema}

\begin{problema}[la advertencia de MATLAB]\label{prob:un-advertencia}
\etiquetaM{} Comente la opción \cmd{'Events'} en \cmd{C2\_Ej9\_Escape} y ejecute de nuevo. Copie la advertencia y explique, con la solución exacta, por qué aparece y en qué instante.
\end{problema}

\begin{problema}[detecte la trampa]\label{prob:un-trampa}
Un estudiante afirma: ``$\dot x=-x+x^3$ tiene el origen localmente estable, luego toda solución converge al origen''. Encuentre el error, determine la región de atracción del origen y calcule el tiempo de escape para $x_0>1$.
\end{problema}

\begin{problema}[un sistema de primer orden con entrada]\label{prob:un-entrada}
\etiquetaEstrella{} Para $\dot x=-x^3+u$ con $u$ constante, encuentre el equilibrio en función de $u$ y su estabilidad. Si $u(t)=\sin t$, ¿puede oscilar la solución? ¿Contradice eso la \cref{prop:un-monotonia}?
\end{problema}

%% ---------------------------------------------------------------------
\section{Soluciones seleccionadas}
%% ---------------------------------------------------------------------

\paragraph{Problema~\ref{prob:un-sillanodo}}
$r<0$: dos equilibrios, $-\sqrt{-r}$ estable y $\sqrt{-r}$ inestable. $r=0$: uno, semiestable. $r>0$: ninguno; todas las soluciones escapan en tiempo finito. En el estable, $f'=-2\sqrt{-r}\to0$: la escala de tiempo $1/\abs{f'}$ crece sin límite y la convergencia se vuelve lenta (es el retraso de la \cref{tr:un-histeresis}).

\paragraph{Problema~\ref{prob:un-allee}}
$0$ estable, $a$ inestable, $1$ estable. Región de atracción de $0$: $(-\infty,a)$ (en el modelo, $[0,a)$); de $1$: $(a,\infty)$. Como en la cosecha, el equilibrio inestable es un umbral de extinción; la diferencia es que aquí el umbral lo pone la biología (encontrar pareja), no el cosechador.

\paragraph{Problema~\ref{prob:un-finito}}
$\dif x/x^{1/3}=-\dif t$ da $\tfrac32\bigl(x_0^{2/3}-x^{2/3}\bigr)=t$, que llega a $0$ en $T$. No contradice la proposición porque $x^{1/3}$ no es Lipschitz en $0$. Si $f$ es localmente Lipschitz y $x(T)=\eq$, entonces $x(t)$ y la constante $\eq$ resuelven el mismo problema de valor inicial en $T$, hacia atrás en el tiempo; por unicidad, $x\equiv\eq$, contradicción.

\paragraph{Problema~\ref{prob:un-clasificar}}
(a) $0$ estable, $\pm1$ inestables. (b) $k\pi$; $f'(k\pi)=e^{-k\pi}(-1)^k$: alternan como en $\sin x$. (c) $0$ con $f'(0)=0$: semiestable ($f\ge0$ a los dos lados); $1$ estable. (d) $0$, con $f'(0)=-1+2=1$: inestable; y otra raíz positiva cerca de $0.797$, estable. La proposición falla en (c), y se decide por el signo de $f$.

\paragraph{Problema~\ref{prob:un-coordinacion}}
$\dot p_1=p_1(1-p_1)\bigl((a+b)p_1-b\bigr)$: equilibrio interior $b/(a+b)$ inestable; $0$ y $1$ estables, con regiones $[0,b/(a+b))$ y $(b/(a+b),1]$. Biestable, como el diodo túnel: la población ``recuerda'' qué estrategia era mayoritaria.

\paragraph{Problema~\ref{prob:un-trampa}}
El error es pasar de local a global. Los equilibrios son $0$ (estable) y $\pm1$ (inestables); la región de atracción del origen es $(-1,1)$. Para $x_0>1$, $\bar t=\tfrac12\ln\frac{x_0^2}{x_0^2-1}$.

\respuestaspausas

